\section{Comparative Synthesis: What Looks Promising?}
The corpus supports four nested levels of decision. At the material level, surface chemistry, roughness, hydration, lubrication, and active particles can alter attachment or removal. At the biological level, effects depend on fouling stage and organism. At the engineering level, roughness, wear, cleaning, and motion determine retained function. At the sustainability level, release, ecotoxicity, energy, cost, and disposal determine whether the intervention is preferable overall.

\subsection{Quality-weighted interpretation}
Study quality was assessed with the four-component Study Reliability Index (SRI) introduced in Section II-C, covering duration, realism, hydrodynamic characterization, and reporting completeness. The rubric and re-tabulated conclusions are provided in Supplementary Table S3. The high-SRI tier is a minority of the corpus, so higher-weight conclusions are reserved for field, operating-asset, or well-characterized dynamic evidence; short laboratory and material-property studies remain mechanistic or exploratory evidence. This changes the interpretation without discarding useful studies: nanocomposite activity and low-adhesion behavior are supported mainly at intermediate SRI, whereas retained performance, release, and ship-scale benefit remain high-SRI evidence gaps for many systems.

No single technology dominates all four levels. Silicone and lubricated interfaces are attractive for low adhesion but require retention and cleaning compatibility. Nanocomposites can add activity or strength but introduce dispersion and release questions. Active UV, photocatalytic, ultrasonic, or laser systems can supplement passive control but require energy, access, and safety. Mechanical grooming and targeted cleaning remain necessary because no coating prevents all fouling under all exposure conditions. The most defensible architecture is a condition-based combination of low-adhesion coating, limited active function where justified, monitoring, and controlled cleaning.

\section{Research Gaps and Future Research Agenda}
\subsection{Standardized, multi-stage validation}
Studies should report substrate preparation, coating thickness, roughness, chemistry, loading, organism or community, flow, exposure, cleaning history, endpoint, and uncertainty. Static settlement, adhesion strength, release, and hydrodynamic drag should be linked rather than reported as isolated metrics.

\subsection{Ageing and full-scale validation}
The largest deployment gap is combined ageing under immersion, abrasion, fouling, UV, pressure, and cleaning. Long-duration field trials should pair surface characterization with fouling, release, resistance, fuel or power, and maintenance records.

\subsection{Environmental and economic integration}
Future work should quantify coating and nanomaterial release, cleaning wastewater, non-target effects, replacement, and disposal. Life-cycle cost and environmental assessment should use the same functional unit and operating scenario so that a lower-fouling claim can be tested against the burden of manufacturing and maintenance.

\subsection{Reproducible benchmarking}
An evidence repository should preserve raw endpoint definitions, negative results, failed formulations, and NR fields. Comparable benchmark panels across laboratories and assets would make it possible to separate chemistry effects from test-condition effects.
